% =============================================================================
%  Кошторис Meet — обсяг робіт, терміни, витрати
%
%  Компілювати XeLaTeX:
%      tectonic docs/koshtorys.tex
%      # або:  xelatex docs/koshtorys.tex   (двічі, заради змісту)
%
%  Шрифти беруться з дистрибутиву TeX (Noto Sans), а не з системи,
%  тож файл однаково збирається на macOS, Linux та в Overleaf.
% =============================================================================

\documentclass[11pt,a4paper]{article}

\usepackage{fontspec}
\usepackage{polyglossia}
\setdefaultlanguage{ukrainian}

\setmainfont{NotoSans-Regular.ttf}[
  BoldFont       = NotoSans-Bold.ttf,
  ItalicFont     = NotoSans-Italic.ttf,
  BoldItalicFont = NotoSans-BoldItalic.ttf,
  Ligatures      = TeX,
]
\setmonofont{NotoSansMono-Regular.ttf}[Scale=0.86]

\usepackage[a4paper,top=22mm,bottom=22mm,left=20mm,right=20mm]{geometry}
\usepackage[table]{xcolor}
\usepackage{booktabs}
\usepackage{tabularx}
\usepackage{array}
\usepackage{ragged2e}
\usepackage{enumitem}
\usepackage{titlesec}
\usepackage[most]{tcolorbox}
\usepackage{fancyhdr}
\usepackage{hyperref}
\usepackage{longtable}
\usepackage{xltabular}

% ---------- палітра (та сама, що й у веб-версії) ----------
\definecolor{ink}{HTML}{121A18}
\definecolor{ink2}{HTML}{33413D}
\definecolor{muted}{HTML}{64716D}
\definecolor{line}{HTML}{CDD5D1}
\definecolor{linesoft}{HTML}{DDE3E0}
\definecolor{surface}{HTML}{F5F7F6}
\definecolor{accent}{HTML}{2C6A5B}
\definecolor{accentsoft}{HTML}{E4EEEA}
\definecolor{critical}{HTML}{9E2A2A}
\definecolor{criticalsoft}{HTML}{F5E4E1}
\definecolor{warn}{HTML}{806011}

\hypersetup{
  colorlinks = true, linkcolor = accent, urlcolor = accent,
  pdftitle   = {Кошторис Meet},
  pdfsubject = {Обсяг робіт, терміни та витрати до релізу},
}

\color{ink2}

% ---------- заголовки ----------
\titleformat{\section}
  {\color{ink}\bfseries\Large}{\thesection.}{0.5em}{}
\titlespacing*{\section}{0pt}{18pt}{6pt}

\titleformat{\subsection}
  {\color{ink}\bfseries\large}{}{0em}{}
\titlespacing*{\subsection}{0pt}{14pt}{4pt}

\titleformat{\subsubsection}
  {\color{ink}\bfseries\normalsize}{}{0em}{}
\titlespacing*{\subsubsection}{0pt}{10pt}{3pt}

% ---------- колонтитули ----------
\pagestyle{fancy}
\fancyhf{}
\renewcommand{\headrulewidth}{0pt}
\renewcommand{\footrulewidth}{0.4pt}
\renewcommand{\footrule}{\color{line}\hrule height 0.4pt\vspace{4pt}}
\fancyfoot[L]{\color{muted}\footnotesize Кошторис Meet · 17 серпня 2026}
\fancyfoot[R]{\color{muted}\footnotesize \thepage}

% ---------- таблиці ----------
\renewcommand{\arraystretch}{1.35}
\arrayrulecolor{linesoft}
\setlength{\tabcolsep}{6pt}

\newcolumntype{L}[1]{>{\RaggedRight\arraybackslash}p{#1}}
\newcolumntype{R}[1]{>{\RaggedLeft\arraybackslash}p{#1}}
\newcolumntype{Y}{>{\RaggedRight\arraybackslash}X}
% перша колонка таблиць: без переносів, інакше короткі назви ріжуться навпіл
\newcolumntype{N}[1]{>{\RaggedRight\arraybackslash\hyphenpenalty=10000\exhyphenpenalty=10000}p{#1}}

% шапка таблиці
\newcommand{\thead}[1]{{\color{muted}\footnotesize\bfseries #1}}

% ---------- позначки типу роботи ----------
\newcommand{\must}{\colorbox{criticalsoft}{\color{critical}\scriptsize\bfseries\,ТРЕБА\,}}
\newcommand{\rec}{\colorbox{accentsoft}{\color{accent}\scriptsize\bfseries\,ВАРТО\,}}

% ---------- врізки ----------
\newtcolorbox{note}[1][]{
  enhanced, breakable, colback=surface, colframe=line,
  boxrule=0.4pt, leftrule=2pt, borderline west={2pt}{0pt}{critical},
  left=10pt, right=10pt, top=8pt, bottom=8pt, arc=0pt, outer arc=0pt,
  fontupper=\normalsize, #1
}
\newtcolorbox{notepos}[1][]{
  enhanced, breakable, colback=surface, colframe=line,
  boxrule=0.4pt, borderline west={2pt}{0pt}{accent},
  left=10pt, right=10pt, top=8pt, bottom=8pt, arc=0pt, outer arc=0pt,
  fontupper=\normalsize, #1
}
\newcommand{\notelbl}[2]{{\color{#1}\footnotesize\bfseries\MakeUppercase{#2}}\par\vspace{3pt}}

\setlist[itemize]{leftmargin=1.2em, itemsep=2pt, topsep=3pt}

\begin{document}

% =============================================================================
%  Титул
% =============================================================================
\begin{center}
  {\color{muted}\footnotesize\bfseries ОБСЯГ РОБІТ · ТЕРМІНИ · ВИТРАТИ}\\[10pt]
  {\color{ink}\fontsize{30}{34}\selectfont\bfseries Кошторис Meet}\\[12pt]
  \begin{minipage}{0.86\textwidth}
    \centering
    Що залишається зробити, щоб додаток вийшов у App~Store і Google~Play:
    обсяг за напрямами, оцінка в робочих днях і повний перелік витрат на
    сервіси --- разово та щомісяця.
  \end{minipage}\\[14pt]
  {\color{muted}\footnotesize 17 серпня 2026 \quad·\quad Оцінка для одного розробника \quad·\quad Ціни в USD}
\end{center}

\vspace{6pt}
{\color{ink}\hrule height 1.2pt}
\vspace{12pt}

% ---------- підсумкові показники ----------
\begin{center}
\setlength{\tabcolsep}{10pt}
\begin{tabular}{ccccc}
  {\color{ink}\Large\bfseries 23 дн} &
  {\color{accent}\Large\bfseries 37 дн} &
  {\color{ink}\Large\bfseries 7--9} &
  {\color{ink}\Large\bfseries \$136} &
  {\color{ink}\Large\bfseries \$25+} \\[2pt]
  {\color{muted}\scriptsize мінімум} &
  {\color{muted}\scriptsize рекомендовано} &
  {\color{muted}\scriptsize тижнів} &
  {\color{muted}\scriptsize разово} &
  {\color{muted}\scriptsize на місяць} \\
\end{tabular}
\end{center}

\vspace{10pt}
{\color{line}\hrule height 0.4pt}

% =============================================================================
\section{Що входить у кошторис}

Документ описує \textbf{лише те, що залишається зробити}. Роботи з бекенду,
безпеки даних, продуктивності та алгоритму підбору вже виконані й у кошторис
не входять.

Оцінки --- у робочих днях однієї людини на повній зайнятості. Календарні
терміни довші за суму днів, бо частина процесів чекає на зовнішні сторони:
перевірку в магазинах і обов'язковий період тестування.

\begin{note}
\notelbl{critical}{Головне обмеження в графіку}
Google вимагає від нових особистих акаунтів розробника \textbf{12 тестерів
протягом 14 днів поспіль} до першого публічного релізу. Це календарний час,
який не можна прискорити грошима чи людьми --- тестування треба запускати на
початку останнього етапу, а не в кінці.
\end{note}

% =============================================================================
\section{Обсяг за напрямами}

\subsection{Інтерфейс і досвід користувача \quad {\color{accent}12 днів}}

\begin{xltabular}{\textwidth}{N{3.7cm} Y R{1.1cm} L{1.5cm}}
\toprule
\thead{Робота} & \thead{Що змінюється і навіщо} & \thead{Днів} & \thead{Тип} \\
\midrule

\textbf{Онбординг} &
Обов'язковий крок після реєстрації: фото, інтереси, дозвіл на геолокацію.
Зараз людина потрапляє одразу в стрічку без жодних даних про себе --- а отже
отримує найгірші можливі рекомендації рівно в той момент, коли вирішує,
залишатись чи ні. & 3 & \must \\

\textbf{Єдина тема, темний режим} &
Кольори зараз прописані в кожному екрані окремо. Це і причина, чому темний
режим виглядає зламано на сучасних телефонах, і причина, чому будь-яка зміна
оформлення коштує дорого. & 3 & \rec \\

\textbf{Локалізація} &
68 рядків жорстко вписані українською просто в екрани, включно з текстами
сповіщень. Іспанська, польська та португальська недоперекладені ---
іноземець бачить наполовину українську програму. & 3 & \rec \\

\textbf{Порожні стани та завантаження} &
Skeleton замість спінера дає відчуття вдвічі швидшого застосунку без жодної
оптимізації. Кожен порожній екран має пропонувати наступний крок, а не
повідомляти, що нічого немає. & 2 & \rec \\

\textbf{Розділи-заглушки} &
«Сповіщення» і «Приватність» у налаштуваннях --- текст-заповнювач. Рев'юер
магазину це відкриє. & 1 & \must \\

\bottomrule
\end{xltabular}

\subsection{Довіра, модерація, безпека людей \quad {\color{accent}13 днів}}

Це не побажання, а формальна вимога Apple до будь-якого застосунку з
користувацьким контентом. Без неї подача буде відхилена.

\begin{xltabular}{\textwidth}{N{3.7cm} Y R{1.1cm} L{1.5cm}}
\toprule
\thead{Робота} & \thead{Що змінюється і навіщо} & \thead{Днів} & \thead{Тип} \\
\midrule

\textbf{Скарги і блокування} &
Зараз кнопка «Заблокувати» просто відхиляє заявку --- реального блокування не
існує. Потрібні скарги, блокування та фільтрація заблокованих з усіх списків
і стрічок. & 3 & \must \\

\textbf{Панель модерації} &
Проста внутрішня сторінка для розгляду скарг. Apple вимагає реакції протягом
24 годин, а без інструмента це фізично неможливо. & 2 & \must \\

\textbf{Верифікація профілю} &
Селфі з бейджем у профілі. Найсильніший захист від ботів і фейків, головний
аргумент довіри для користувачів і вагомий плюс при перевірці магазином. & 3 & \must \\

\textbf{Вікове обмеження} &
Дата народження збирається, але перевірки на 18+ немає. Для продукту з
офлайн-зустрічами це критично і поза магазинами. & 1 & \must \\

\textbf{Безпека зустрічей} &
Поділитись локацією з другом, тривожна кнопка, підказки про публічні місця.
Від застосунку, що зводить незнайомців офлайн, це прямо очікують --- і
користувачі, і рев'юери. & 4 & \rec \\

\bottomrule
\end{xltabular}

\subsection{Технічна частина \quad {\color{accent}5 днів}}

\begin{xltabular}{\textwidth}{N{3.7cm} Y R{1.1cm} L{1.5cm}}
\toprule
\thead{Робота} & \thead{Що змінюється і навіщо} & \thead{Днів} & \thead{Тип} \\
\midrule

\textbf{Захист фотографій} &
Зараз фото доступне за прямим посиланням тому, хто його має. Перехід на
посилання з обмеженим строком дії закриває це остаточно. & 2 & \rec \\

\textbf{Підвантаження стрічки} &
Зараз стрічка обмежена однією порцією: дійшов до кінця --- і все до
перезапуску застосунку. Для продукту, що живе з переглядів, це прямо ріже
утримання. & 1,5 & \must \\

\textbf{Моніторинг помилок} &
Підключення та перевірка збору збоїв. Механізм готовий, потрібен обліковий
запис і налаштування. & 0,5 & \must \\

\textbf{Регламентні завдання} &
Автоматичне оновлення ваг інтересів і очищення службових таблиць за
розкладом. & 1 & \rec \\

\bottomrule
\end{xltabular}

\subsection{Випуск у магазини \quad {\color{accent}7 днів}}

\begin{xltabular}{\textwidth}{N{3.7cm} Y R{1.1cm} L{1.5cm}}
\toprule
\thead{Робота} & \thead{Що змінюється і навіщо} & \thead{Днів} & \thead{Тип} \\
\midrule

\textbf{Айдентика застосунку} &
Назва, іконка, заставка, ідентифікатор. Зараз усе тимчасове: застосунок
називається «Dating App». & 2 & \must \\

\textbf{Підпис і дозволи} &
Ключ підпису релізу --- зараз збірка підписана тимчасовим ключем і Google її
не прийме. Плюс тексти запитів дозволів для iOS, без яких застосунок падає
при виборі фото. & 1 & \must \\

\textbf{Юридичні сторінки} &
Політика конфіденційності, умови використання, публічна сторінка видалення
акаунта. Усі три --- обов'язкові посилання при подачі. & 1 & \must \\

\textbf{Декларації про дані} &
Анкети про збір даних в обох магазинах: локація, фото, телефон,
ідентифікатори. Помилка тут --- відхилення або зняття з публікації пізніше. & 1 & \must \\

\textbf{Матеріали і подача} &
Скріншоти під потрібні розміри екранів, описи, збірки, завантаження,
супровід перевірки. & 2 & \must \\

\bottomrule
\end{xltabular}

% =============================================================================
\section{Два обсяги на вибір}

\vspace{4pt}
\noindent
\begin{minipage}[t]{0.475\textwidth}
\begin{tcolorbox}[enhanced, colback=surface, colframe=line, boxrule=0.4pt,
                  arc=0pt, outer arc=0pt, left=10pt, right=10pt, top=9pt, bottom=9pt]
{\color{muted}\footnotesize\bfseries ВАРІАНТ A · МІНІМУМ}\\[6pt]
{\color{ink}\LARGE\bfseries 23 дні}\\[2pt]
{\color{muted}\footnotesize $\approx$ 5 робочих тижнів · подача без відхилення}\\[8pt]
\begin{itemize}[leftmargin=1em]
  \item Уся модерація і верифікація
  \item Онбординг
  \item Підвантаження стрічки
  \item Заміна розділів-заглушок
  \item Захист фотографій, моніторинг
  \item Повний випуск у магазини
\end{itemize}
\vspace{4pt}
{\color{muted}\footnotesize Застосунок пройде перевірку і працюватиме. Без темного
режиму, без повної локалізації, без функцій безпеки зустрічей.}
\end{tcolorbox}
\end{minipage}
\hfill
\begin{minipage}[t]{0.475\textwidth}
\begin{tcolorbox}[enhanced, colback=surface, colframe=line, boxrule=0.4pt,
                  arc=0pt, outer arc=0pt, left=10pt, right=10pt, top=9pt, bottom=9pt,
                  borderline north={2pt}{0pt}{accent}]
{\color{accent}\footnotesize\bfseries ВАРІАНТ B · РЕКОМЕНДОВАНИЙ}\\[6pt]
{\color{ink}\LARGE\bfseries 37 днів}\\[2pt]
{\color{muted}\footnotesize $\approx$ 7,5 робочих тижнів · продукт, а не мінімум}\\[8pt]
\begin{itemize}[leftmargin=1em]
  \item Усе з варіанта A
  \item Єдина тема і темний режим
  \item Повна локалізація 4 мовами
  \item Порожні стани і skeleton
  \item Безпека зустрічей: SOS, локація
  \item Регламентні завдання
\end{itemize}
\vspace{4pt}
{\color{muted}\footnotesize Різниця в 14 днів --- це те, що відрізняє застосунок,
який пропустили, від застосунку, яким користуються.}
\end{tcolorbox}
\end{minipage}

\vspace{10pt}

\begin{notepos}
\notelbl{accent}{Якщо треба скоротити}
Найдешевше відкласти \textbf{темний режим} і \textbf{локалізацію} --- це 6 днів,
і обидва пункти безболісно робляться після запуску. \textbf{Онбординг
скорочувати не варто ні за яких умов}: це три дні роботи, які найсильніше
впливають на те, чи залишиться користувач після першого відкриття.
\end{notepos}

% =============================================================================
\section{Календарний план}

За рекомендованим обсягом, з урахуванням очікування на зовнішні сторони.

\vspace{6pt}
\begin{xltabular}{\textwidth}{N{5.2cm} Y L{2.8cm}}
\toprule
\thead{Етап} & \thead{Характер} & \thead{Період} \\
\midrule
Модерація і довіра        & робота                                & тижні 1--2 \\
Інтерфейс                 & робота                                & тижні 3--4 \\
Технічна частина          & робота                                & тиждень 5 \\
Айдентика і збірки        & робота                                & тижні 5--6 \\
Тестування Google         & {\color{warn}очікування, 14 днів поспіль} & тижні 6--8 \\
Перевірка магазинами      & {\color{warn}очікування}              & тиждень 9 \\
\bottomrule
\end{xltabular}

\vspace{6pt}
Обов'язкові 14 днів тестування Google запускаються паралельно зі збірками,
тому загальний строк --- \textbf{7--9 тижнів}, а не сума всіх етапів.

% =============================================================================
\section{Витрати на сервіси}

\subsection{Разові платежі}

\begin{xltabular}{\textwidth}{N{5cm} Y R{2.6cm}}
\toprule
\thead{Стаття} & \thead{Коментар} & \thead{Вартість} \\
\midrule
\textbf{Google Play Console}     & Одноразово, назавжди & \$25 \\
\textbf{Apple Developer Program} & Щорічно; без нього публікація в App Store неможлива & \$99 / рік \\
\textbf{Домен}                   & Для юридичних сторінок і сторінки видалення акаунта & $\sim$\$12 / рік \\
\midrule
{\color{ink}\bfseries Разом на старті} & & {\color{ink}\bfseries \$136} \\
\bottomrule
\end{xltabular}

\subsection{Щомісячні витрати за масштабом}

Три сценарії: закрита бета, перше місто, зростання.

\begin{xltabular}{\textwidth}{Y R{2.5cm} R{2.5cm} R{2.7cm}}
\toprule
\thead{Стаття} & \thead{100 активних} & \thead{1 000 активних} & \thead{10 000 активних} \\
\midrule
\textbf{База даних і сховище} \newline {\color{muted}\footnotesize Supabase}
  & \$0 & \$25 & \$25 + трафік \\
\textbf{Трафік на фото} \newline {\color{muted}\footnotesize понад включений обсяг}
  & \$0 & \$0--20 & \$80--250 \\
\textbf{SMS-підтвердження} \newline {\color{muted}\footnotesize головна змінна стаття}
  & \$5--10 & \$50--90 & \$300--700 \\
\textbf{Моніторинг помилок} \newline {\color{muted}\footnotesize Sentry}
  & \$0 & \$0--26 & \$26 \\
\textbf{Сповіщення} \newline {\color{muted}\footnotesize Firebase}
  & \$0 & \$0 & \$0 \\
\midrule
{\color{ink}\bfseries На місяць}
  & {\color{ink}\bfseries \$5--10} & {\color{ink}\bfseries \$75--160} & {\color{ink}\bfseries \$430--1000} \\
\bottomrule
\end{xltabular}

\vspace{4pt}
{\color{muted}\footnotesize Ціни зовнішніх сервісів наведено орієнтовно станом на
дату документа і їх варто звірити при оформленні. «Активний користувач» тут ---
людина, яка відкриває застосунок кілька разів на місяць.}

\vspace{8pt}
\begin{note}
\notelbl{critical}{Найбільший ризик у витратах}
\textbf{SMS домінують над усім іншим разом узятим} і зростають лінійно з
реєстраціями. Один код підтвердження коштує помітних грошей, а вартість SMS
на Україну --- одна з найвищих у Європі.

\vspace{4pt}
Рішення просте й окупається одразу: додати \textbf{вхід через Google та Apple}
і залишити SMS запасним варіантом. Більшість людей обирає вхід у два дотики,
і стаття витрат скорочується в рази. Робота --- \textbf{2 дні}, і вона
економить більше, ніж усі інші оптимізації разом.

\vspace{4pt}
Вхід через Apple, до речі, і так стане обов'язковим, щойно в застосунку
з'явиться будь-який інший сторонній вхід --- це вимога App~Store.
\end{note}

\subsection{Що не потребує оплати}

Сповіщення, зберігання коду, сторінки політики конфіденційності та видалення
акаунта, середовище розробки, збірка застосунків. Усе це покривається
безкоштовними тарифами й не з'явиться в рахунках навіть при зростанні.

\subsection{Чого немає в кошторисі}

\begin{itemize}
  \item \textbf{Просування} --- залучення перших користувачів планується окремо
  \item \textbf{Юридичний супровід} --- політику конфіденційності можна зробити
        за шаблоном, але перевірка юристом коштуватиме окремо
  \item \textbf{Дизайнер} --- якщо іконку та фірмовий стиль замовляти на стороні,
        це \$100--400 одноразово
  \item \textbf{Розумніший підбір за описом профілю} --- можлива наступна
        ітерація, орієнтовно 4 дні роботи і одиниці доларів на місяць
\end{itemize}

% =============================================================================
\section{Припущення і ризики}

\begin{xltabular}{\textwidth}{N{5.4cm} Y}
\toprule
\thead{Припущення} & \thead{Що станеться, якщо не справдиться} \\
\midrule
\textbf{Один розробник на повній зайнятості} &
При частковій зайнятості календарний строк зростає пропорційно;
сума днів не змінюється. \\

\textbf{Перевірка з першого разу} &
Відхилення App Store додає 3--7 днів на виправлення й повторну подачу.
Перша подача проходить із першого разу приблизно в половині випадків. \\

\textbf{Юридичні тексти за шаблоном} &
Якщо потрібен супровід юриста --- плюс кілька днів очікування і окремий
рахунок. \\

\textbf{Запуск в одному місті} &
Витрати на SMS і трафік рахувалися для концентрованої аудиторії. Ширший
запуск підвищує витрати, не підвищуючи щільності --- а саме щільність
визначає, чи працює продукт. \\
\bottomrule
\end{xltabular}

\vspace{12pt}

\begin{notepos}
\notelbl{accent}{Підсумок}
\textbf{Рекомендований обсяг:} 37 робочих днів, 7--9 тижнів календарно.\\
\textbf{Витрати на старті:} \$136 разово, від \$25 на місяць.\\
\textbf{При тисячі активних користувачів:} \$75--160 на місяць, з яких
більшість --- SMS.\\
\textbf{Найкраща інвестиція в кошторисі:} 2 дні на вхід через Google та Apple ---
вони окуповуються з першої сотні реєстрацій.
\end{notepos}

\end{document}
